% !TEX TS-program = lualatex
% !TEX encoding = UTF-8

\documentclass[12pt, twoside, french, openany]{book}
\input{headers}
\externaldocument[M-]{\subfix{VesperaleOP}}

\begin{document}

%\customsubfile{vop01_prooemium}
\customsubfile{vop03_tempus_adventus}
\customsubfile{vop05_tempus_nativitatis}
\customsubfile{vop06_tempus_post_epiphaniam}
\customsubfile{vop07_tempus_quadragesimae}
\customsubfile{vop09_tempus_paschale}
\customsubfile{vop10_tempus_pentecosten}
%\customsubfile{vop11_tempus_post_trinitatem}
\customsubfile{vop13_ordinarium}
\customsubfile{vop15_psalterium}
\customsubfile{vop17_psalterium_tp}
%\customsubfile{vop21_commune_bmv}
\customsubfile{vop23_commune_dedicationis}
\customsubfile{vop25_commune_apostolorum}
\customsubfile{vop27_commune_martyrum}
\customsubfile{vop28_commune_plurimorum_martyrum}
\customsubfile{vop29_commune_confessorum}
\customsubfile{vop30_commune_virginum_martyrum}
\customsubfile{vop31_commune_virginum_non_martyrum}
\customsubfile{vop32_commune_non_virginum}
\customsubfile{vop33_communia_tp}
%\customsubfile{vop37_officium_defunctorum}
%\customsubfile{vop40_festa_mobilia}
\customsubfile{vop50_sanctorale_decembris}
\customsubfile{vop51_sanctorale_januarii}
\customsubfile{vop52_sanctorale_februarii}
%\customsubfile{vop53_sanctorale_martii}
%\customsubfile{vop54_sanctorale_aprilis}
%\customsubfile{vop55_sanctorale_maii}
%\customsubfile{vop56_sanctorale_junii}
%\customsubfile{vop57_sanctorale_julii}
%\customsubfile{vop58_sanctorale_augusti}
%\customsubfile{vop59_sanctorale_septembris}
%\customsubfile{vop60_sanctorale_octobris}
%\customsubfile{vop61_sanctorale_novembris}

%% définition du style de la section d'index
\cleardoublepage
\addcontentsline{toc}{chapter}{Index}
\setheader{Index}
\pagestyle{fancy}
\thispagestyle{empty}
%% impression des index
\cprintindex{A}{Index Antiphonarum} % antiennes
\cprintindex{H}{Index Hymnorum} % hymnes
\cprintindex{R}{Index Responsoriorum} % répons
%\cprintindex{F}{Index Festorum} % fêtes. Actuellement cassé.

\cleardoublepage % on avance jusqu'à la première page de droite
%% On donne son titre à la table des matières principale
\renewcommand{\contentsname}{{\centering\normalfont\scshape Index generalis\par}}
\thispagestyle{empty}
\tableofcontents % on imprime la table des matières

\end{document}
